\section{Problem setup}
\label{sec:problem-setup}

We study aggregate-supervised bags of instances.
A training example is a bag $B=(x_1,\ldots,x_n)$ together with an observed aggregate label $y$.
The instance-level variables that explain $y$ are hidden during training and are used only for evaluation.
The aggregate label may be a Bernoulli count, a signed count, an integer-valued sum, or a class histogram.
In all cases, each instance contributes a finite-support discrete value, and the observed label is the sum of these contributions.
We define \method{} as one likelihood layer for this structure, so different aggregate-supervision tasks require only different atomic supports and contribution maps.

The unifying object is the \emph{atomic support}.
For a one-dimensional task, choose a finite set
\begin{equation}
  \mathcal{A}=\{v_0,\ldots,v_{S-1}\}\subset\mathbb{Z},
\end{equation}
where $S=|\mathcal{A}|$ is the atomic support size.
A backbone model produces scores for each instance, and a contribution map turns those scores into an atomic PMF over $\mathcal{A}$.
Changing $\mathcal{A}$ changes the aggregate-supervision task while leaving the aggregate-prediction algorithm unchanged:
$S=2$ gives Bernoulli counts, $S=10$ gives digit sums over $\{0,\ldots,9\}$, and $S=3$ gives signed atoms over $\{-1,0,+1\}$.
For bag size $n$ and contiguous support $\{0,\ldots,S-1\}$, the aggregate PMF has length
\begin{equation}
  T = n(S-1)+1.
  \label{eq:aggregate-support-width}
\end{equation}
Signed supports use the same principle with an offset, and more general atoms can use instance-dependent supports.

Let $C_i$ denote the latent contribution of instance $x_i$.
The contribution may also depend on observed side information $a_i$, such as a known sign in signed-count experiments.
We assume
\begin{equation}
  C_i \in \mathcal{S}_i \subset \mathbb{Z}^d,
  \qquad
  |\mathcal{S}_i| < \infty,
\end{equation}
and the observed aggregate is
\begin{equation}
  Y = \sum_{i=1}^n C_i.
  \label{eq:aggregate-sum}
\end{equation}
Most experiments use $d=1$.
For histograms, $d$ may be the number of classes, although our scalable implementation uses one-dimensional one-vs-rest count likelihoods.
The learner observes $(x_{1:n},a_{1:n},y)$ when side information is present, but never observes the individual $C_i$ values during training.
When discussing individual scalar tasks, we write $Y_{\rm count}$ for ordinary nonnegative counts, $Y_{\rm sum}$ for integer-valued sums such as digit sums, and $Y_{\pm}$ for signed counts.
These are task-specific instances of the generic aggregate $Y$ in Eq.~\eqref{eq:aggregate-sum}.

\begin{definition}[\method{} finite-support convolutional likelihood layer]
\label{def:finite-support-aggregator}
Let $f_\theta$ be any instance model, and let $\Phi$ be a contribution map that converts the output of $f_\theta$ on instance $x_i$ and observed side information $a_i$ into an atomic PMF
\begin{equation}
  d_i = \Phi(f_\theta(x_i),a_i) \in \Delta(\mathcal{S}_i),
  \qquad
  \mathcal{S}_i \subset \mathbb{Z}^d,\quad |\mathcal{S}_i|<\infty .
\end{equation}
For conditionally independent latent contributions $C_i\sim d_i$, \method{} is the deterministic map
\begin{equation}
  \mathsf{Agg}(d_1,\ldots,d_n)(y)
  =
  (d_1 * \cdots * d_n)(y)
  =
  \sum_{\substack{c_1\in\mathcal{S}_1,\ldots,c_n\in\mathcal{S}_n\\
                  c_1+\cdots+c_n=y}}
  \prod_{i=1}^n d_i(c_i),
  \label{eq:finite-support-aggregator}
\end{equation}
which assigns the exact probability mass of the aggregate $Y=\sum_i C_i$ to value $y$.
\end{definition}

For Bernoulli atoms $d_i=[1-p_i,p_i]$ on $\{0,1\}$, \method{} recovers the count-loss probability
\begin{equation}
  \Pr_\theta(Y=k)=(d_1*\cdots*d_n)(k)=
  [t^k]\prod_{i=1}^n \bigl((1-p_i)+p_it\bigr),
\end{equation}
the Poisson-binomial probability optimized by Bernoulli count-loss methods.
Under identical instance probabilities, \method{} and CountLoss define the same loss and gradients up to backend differences.
Applying the same equation independently to each class with atoms $[1-p_{ic},p_{ic}]$ recovers the one-vs-rest count-likelihood component used in LLP-PVC-style objectives, but not a normalized full joint multiclass histogram likelihood.

Definition~\ref{def:finite-support-aggregator} separates four choices often coupled in task-specific aggregate-supervision losses:
\begin{itemize}
  \item \emph{Backbone independence.}  The network supplies atomic PMFs; the likelihood layer is unchanged across MLPs, CNNs, transformers, or pretrained encoders.
  \item \emph{Contribution-map independence.}  Different aggregate-supervision tasks are specified by the contribution map and finite support, not by changing the aggregation rule.
  \item \emph{Backend independence.}  The returned PMF supports NLL training, while its computation can be implemented by dynamic programming, direct convolution, grouped convolution, or FFT-tree multiplication.
  \item \emph{Conditional-independence scope.}  The PMF is exact under conditionally independent latent contributions; otherwise it is a structured aggregate-likelihood approximation.
\end{itemize}
We train with the aggregate negative log likelihood $-\log \mathsf{Agg}(d_1,\ldots,d_n)(y)$ for the observed label $y$.

Once the atomic contribution distributions $d_i$ are predicted, the aggregate-label distribution is fixed by convolution under the conditional-independence model.
Tasks differ in the support $\mathcal{S}_i$, the contribution map $\Phi$, and which side information is observed; concrete kernels are summarized in Table~\ref{tab:method-kernels}.
Bernoulli count supervision is the simplest case of Eq.~\eqref{eq:aggregate-sum}, while signed counts, integer-valued sums, multiplicity counts, and one-vs-rest class counts change the finite-support kernel but keep the same likelihood layer.
In our experiments, bags are sampled on the fly with replacement from the corresponding train or test split.
Thus an image may appear in multiple bags, but train and test bags are always drawn from disjoint dataset splits.
